\documentclass[titlepage,a4paper]{jsarticle}
\usepackage{../../../sty/import}% 各種パッケージインポート
\usepackage{../../../sty/title}% タイトルページの変更

%% タイトルページの変数
% レポートタイトル
\title{量子・原子力統合特論第三回目レポート}
% 提出日
\expdate{\today}
% 科目名
\subject{量子・原子力統合特論}
% 分野
\class{情報経営システム工学分野}
% 学年
\grade{M1}
% 学籍番号
\mynumber{24336488}
% 記述者
\author{本間三暉}

\begin{document}
% titleページ作成
\maketitle
\section{ETIGO-IIで3MeVの陽子ビームをアルミニウムに発射した際に発生しうる放射線の種類を答えよ．}
3 MeVの陽子ビームをアルミニウムに照射した場合，陽子は荷電粒子であるため，まずアルミニウム中の電子とクーロン力によって相互作用し，電離や励起を起こす．
そのため，二次電子，特性X線，オージェ電子が発生しうる．

また，陽子がアルミニウム原子核と相互作用することで核反応が起こる可能性がある．
例えば，
\[{}^{27}\mathrm{Al}(p,\gamma){}^{28}\mathrm{Si}\]や，\[{}^{27}\mathrm{Al}(p,p'\gamma){}^{27}\mathrm{Al}\]
のような反応により，$\gamma$線が発生しうる．
さらに，反応によっては荷電粒子である$\alpha$線が発生する可能性も考えられる．

一方で，中性子を発生する\[{}^{27}\mathrm{Al}(p,n){}^{27}\mathrm{Si}\]のような反応は，しきい値が3 MeVより高いため，3 MeVの陽子ビームでは中性子の発生は主には考えにくい．

したがって，3 MeVの陽子ビームをアルミニウムに照射した際に発生しうる主な放射線は陽子線，二次電子，特性X線，オージェ電子，$\gamma$線，および$\alpha$線である．

\section{ETIGO-IIで3MeVの電子ビームをアルミニウムに発射した際に発生しうる放射線の種類を答えよ．}
3 MeVの電子ビームをアルミニウムに照射した場合，電子は物質中の電子と相互作用し，電離や励起を起こす．これにより，二次電子，特性X線，オージェ電子が発生しうる．

また，高速電子はアルミニウム原子核の電場によって進行方向を曲げられ，減速される．
このとき，電子の運動エネルギーの一部が電磁波として放出されるため，制動X線が発生する．
制動X線の最大エネルギーは入射電子のエネルギーを超えないため，3 MeV電子ビームでは最大約3 MeVのX線が発生しうる．

さらに，発生したX線はアルミニウムや周囲の物質と相互作用し，光電効果やコンプトン散乱を起こす．
そのため，散乱X線や反跳電子も発生しうる．
一方で，アルミニウムから光核反応によって中性子を発生させるには，3 MeVより高いエネルギーの光子が必要であるため，中性子の発生は基本的に考えにくい．

したがって，3 MeVの電子ビームをアルミニウムに照射した際に発生しうる主な放射線は電子線，二次電子，制動X線，特性X線，オージェ電子，散乱X線である．
中性子は主な発生放射線とは考えにくい．

\section{上記の放射線による被爆を防ぐための対策について200字程度で述べよ．}
課題1，2で発生しうる放射線のうち，陽子線や$\alpha$線などの荷電粒子はアクリル，プラスチック，アルミニウム，水，コンクリートなどで遮蔽できる．
電子線や二次電子は，制動X線の発生を抑えるため，まずアクリルやプラスチックなどの低原子番号材料で遮蔽するのがよい．
特性X線，制動X線，散乱X線，$\gamma$線は透過力が大きいため，鉛，鉄，コンクリートなどの高密度材料を用いて減衰させる．
\begin{thebibliography}{99}
\bibitem{}量子原子力統合工学概論 講義スライド
\bibitem{}原子力材料と核燃料 講義スライド
\end{thebibliography}

\end{document}

% ファイル形式は<学籍番号>.pdf